\subsectionpagedleft{Ephemeris}
\label{config:ephemeris}
\index{Ephemeris}
\index{Config!Ephemeris}

\begin{lstlisting}[language=Go,style=kaolstplain,escapechar=|]
ephemeris( "sensor"
    {
        every "crontab" [on start]
        { |\textit{target}|( |\textit{options}\dots| ) [as "name"] }+
    }+
)
\end{lstlisting}

\Statement{ephemeris} calculates an ephemeris of astronomical objects on a specific schedule.

\ParameterDescription{
\item[sensor] is the base sensor Id for the ephemeris.
When the ephemeris is being created, the generated metrics will be formed with the sensor id, the target name followed
by the requested fields.

\item[every]defines when the calculation should be performed.
\marginnote{It's advised to always set \textbf{on start} as it ensures data is available, rather than wait up to
an hour or more before the first calculation is performed.}
For most objects once an hour suffices, although the Sun and Moon are better calculated at shorter periods as they
move fast in relation to the station.
\item[on start]if present then the ephemeris will also be calculated when the \CalculatorService
starts.

\item[target]is the object name to calculate. \vref{tab:config:ephemeris:targets} lists the currently available target names.

\item[options]is a list of fields to include in the ephemeris. \vref{tab:config:ephemeris:options} lists the possible options.
\item[as] overrides the metric generated. If absent then the target name is used.
}

\begin{table}[ht]
 \caption{Available Ephemeris Targets}
 \labtab{config:ephemeris:targets}
 \begin{tabular}{l l l|l}
 \hline
 \multicolumn{3}{l}{\textbf{target}} & \textbf{advised crontab} \\
 \hline
 sun\index{Sun}\index{Ephemeris!Sun} &
     moon\index{Moon}\index{Ephemeris!Moon} &
     & minute \\
 \hline
 mercury\index{Mercury}\index{Ephemeris!Mercury} &
    venus\index{Venus}\index{Ephemeris!Venus} &
    mars\index{Mars}\index{Ephemeris!Mars} &
    \multirow{3}{*}{hour} \\
 jupiter\index{Jupiter}\index{Ephemeris!Jupiter} &
    saturn\index{Saturn}\index{Ephemeris!Saturn} &
    uranus\index{Uranus}\index{Ephemeris!Uranus} & \\
 neptune\index{Neptune}\index{Ephemeris!Neptune} &&& \\
 \hline
 \end{tabular}
\end{table}

\clearpagefit
\begin{table*}[h!]
 \caption{Available Ephemeris Options}
 \labtab{config:ephemeris:options}
 \begin{tabular}{l|l|l|l|l|l}
 \hline
 \multicolumn{4}{c|}{\textbf{option}} & \textbf{purpose} & \textbf{suffix} \\
 \hline
 \multirow{14}{*}{\rotatebox[origin=c]{90}{all}} &
    \multirow{9}{*}{\rotatebox[origin=c]{90}{allCoordinates}} &
        \multirow{2}{*}{equatorial\index{Equatorial}\index{Ephemeris!Equatorial}} & ra & Right Ascension\index{Right Ascension}\index{Ephemeris!Equatorial!Right Ascension}\index{Coordinate!Equatorial!Right Ascension} & eq.ra\\
    &   &   & dec & Declination\index{Declination}\index{Ephemeris!Equatorial!Declination}\index{Coordinate!Equatorial!Declination} & eq.dec \\
        \cline{3-6}
    &   & \multirow{2}{*}{ecliptic\index{Ecliptic}\index{Ephemeris!Ecliptic}} & eclipticLatitude & Ecliptic Latitude\index{Ephemeris!Ecliptic!Latitude}\index{Coordinate!Ecliptic!Latitude} & ecl.lat \\
    &   &   & eclipticLongitude & Ecliptic Longitude\index{Ephemeris!Ecliptic!Longitude}\index{Coordinate!Ecliptic!Longitude} & ecl.lon \\
        \cline{3-6}
    &   & \multirow{2}{*}{galactic\index{Galactic}\index{Ephemeris!Galactic}\index{Cordinate!Galactic}} & galacticLatitude\index{Galactic!Latitude}\index{Ephemeris!Galactic!Latitude}\index{Cordinate!Galactic!Latitude} & Galactic Latitude & gal.lat\\
    &   &   & galacticLongitude & Galactic Longitude\index{Galactic!Longitude}\index{Ephemeris!Galactic!Longitude}\index{Cordinate!Galactic!Longitude} & gal.lon \\
        \cline{3-6}
    &   & \multirow{3}{*}{horizon\index{Horizon}\index{Ephemeris!Horizon}\index{Cordinate!Horizon}} & altitude & Altitude\index{Altitude}\index{Horizon!Altitude}\index{Ephemeris!Horizon!Altitude}\index{Cordinate!Horizon!Altitude}} above horizon & hz.altitude \\
    &   &   & azimuth & Azimuth\index{Azimuth}\index{Horizon!Azimuth}\index{Ephemeris!Horizon!Azimuth}\index{Cordinate!Horizon!Azimuth}} clockwise from \Degrees{0} South & hz.azimuth \\
    &   &   & bearing & Bearing\index{Bearing}\index{Horizon!Bearing}\index{Ephemeris!Horizon!Bearing}\index{Cordinate!Horizon!Bearing}} clockwise from \Degrees{0} North & hz.bearing \\
    \cline{2-6}
 &  \multirow{4}{*}{\rotatebox[origin=c]{90}{allData}} &
            & diameter\index{Diameter}\index{Ephemeris!Diameter} & The objects semi-diameter\index{Semi-diameter} & diameter \\
        \cline{3-6}
    &   & \multirow{3}{*}{distance\index{Distance}\index{Ephemeris!Distance}} & distanceEarth\index{Distance!from Earth}\index{Ephemeris!Distance Earth}} & Distance from the Earth & dist.earth \\
    &   &   & distanceSun\index{Distance!from Sun}\index{Ephemeris!Distance Sun} & Distance from the Sun & dist.sun \\
    &   &   & lightTime\index{Light Time}\index{Ephemeris!Light Time} & Time from Earth at the speed of light & dist.lighttime \\
 \hline
 \end{tabular}
\end{table*}

\begin{widepar}
All options are case-insensitive. Some options include multiple options.
eg equatorial includes ra and dec.
\end{widepar}

\begin{widepar}
Azimuths are of the astronomical kind, so \Degrees{0} is due South not North.
For traditional Azimuth which is clockwise from North then Bearing is used.
\end{widepar}

\subsubsectionfit{Example ephemeris}
\label{subsec:config:ephemeris:simple}
This example shows how to calculate an ephemeris.

\begin{lstlisting}[language=Go,caption={Example ephemeris},escapechar=|]
ephemeris( "sol" |\label{line:ephemeris:simple:sensor}|
    every "minute" on start |\label{line:ephemeris:simple:minute}|
        sun(allData, equatorial, horizon) |\label{line:ephemeris:simple:sun}|
        moon(allData, equatorial, horizon)
    every "hour" on start |\label{line:ephemeris:simple:hour}|
        mercury(allData, equatorial, horizon)
        venus(allData, equatorial, horizon)
        mars(allData, equatorial, horizon)
        jupiter(allData, equatorial, horizon)
        saturn(allData, equatorial, horizon)
        uranus(allData, equatorial, horizon)
        neptune(allData, equatorial, horizon)
    )
\end{lstlisting}

Line \ref{line:ephemeris:simple:sensor} defines the sensor id as "sol".

Line \ref{line:ephemeris:simple:minute} declares we want the following targets calculated every minute and
when the service starts.

Line \ref{line:ephemeris:simple:sun} declares a target. Here it's the Sun and we want to use the options:
allData, equatorial and horizon.
For simplicity in this example, all the targets are using the same options.

Line \ref{line:ephemeris:simple:hour} declares we want the following targets calculated once every hour and
when the service starts.

This example will then generate the following metrics - only the sun and moon are shown here for brevity:

\begin{table}[ht]
    \begin{tabular}{l|l|l|l}
        \hline
        \textbf{option} & \textbf{sun} & \textbf{moon} & \textbf{description} \\
        \hline
        \multirow{2}{*}{equatorial} & sol.sun.eq.ra & sol.moon.eq.ra & Right Ascension \\
            & sol.sun.eq.dec & sol.moon.eq.dec & Declination \\
        \hline
        \multirow{3}{*}{horizon} & sol.sun.hz.altitude & sol.moon.hz.altitude & Altitude \\
            & sol.sun.hz.azimuth & sol.moon.hz.azimuth & Azimuth \\
            & sol.sun.hz.bearing & sol.moon.hz.bearing & Bearing \\
            \hline
        \multirow{4}{*}{allData} & sol.sun.diameter & sol.moon.diameter & Semi-Diameter \\
            & sol.sun.dist.earth & sol.moon.dist.earth & Distance from Earth \\
            & sol.sun.dist.sun & sol.moon.dist.sun & Distance from Sun \\
            & sol.sun.dist.lighttime & sol.moon.dist.lighttime & Light time from Earth \\
            \hline
    \end{tabular}
\end{table}
